% from ega2-1-fr.tex
\oldpage[II]{150}
En effet, le morphisme canonique $f:X'\to X$ est alors fini
(\hyperref[II.7.4.6-fr]{7.4.6}), donc propre
(\hyperref[II.6.1.11-fr]{6.1.11}) et surjectif
(\hyperref[II.6.3.8-fr]{6.3.8})~; si $g:X\to\operatorname{Spec}(k)$ est le
morphisme structural, $g$ et $g\circ f$ sont donc propres simultanément,
comme il résulte de (\hyperref[II.5.4.2-fr]{5.4.2, (iii)}) et
(\hyperref[II.5.4.3-fr]{5.4.3, (ii)}), $g$ étant séparé par hypothèse.

\begin{proposition}[7.4.9]
\label{II.7.4.9-fr}
Soient $X$ une courbe algébrique normale sur $k$, $Y$ un $k$-schéma algébrique
propre sur $k$. Toute $k$-application rationnelle de $X$ dans $Y$ est alors
partout définie, autrement dit est un morphisme.
\end{proposition}

En effet, il résulte de (\hyperref[II.7.3.7-fr]{7.3.7}) qu'aux points $x\in X$
où une telle application n'est pas définie, la dimension de $\mathcal{O}_x$
devrait être $\geq2$, donc l'ensemble de ces points est vide~; la dernière
assertion provient de (\hyperref[I.7.2.3-fr]{I, 7.2.3}).

\begin{corollary}[7.4.10]
\label{II.7.4.10-fr}
Une courbe algébrique normale sur $k$ est quasi-projective sur $k$.
\end{corollary}

Comme $X$ est somme d'un nombre fini de courbes algébriques intègres et
normales (\hyperref[II.6.3.8-fr]{6.3.8}), on peut se borner au cas où $X$ est
intègre (\hyperref[II.5.3.6-fr]{5.3.6}). Comme $X$ est quasi-compact, il est
recouvert par un nombre fini d'ouverts affines $U_i$ $(1\leq i\leq n)$, et
comme chacun de ces derniers est de type fini sur $k$, il existe un entier
$n_i$ et une $k$-immersion $f_i:U_i\to\mathbf{P}_k^{n_i}$
(\hyperref[II.5.3.3-fr]{5.3.3} et
\hyperref[II.5.3.4-fr]{5.3.4, (i)}). Comme $U_i$ est dense dans $X$, il résulte
de (\hyperref[II.7.4.9-fr]{7.4.9}) que $f_i$ se prolonge en un $k$-morphisme
$g_i:X\to\mathbf{P}_k^{n_i}$, d'où un $k$-morphisme
$g=(g_1,\ldots,g_n)_k$ de $X$ dans le produit $P$ des
$\mathbf{P}_k^{n_i}$ sur $k$. En outre, pour chaque indice $i$, comme la
restriction de $g_i$ à $U_i$ est une immersion, il en est de même de la
restriction de $g$ à $U_i$ (\hyperref[I.5.3.14-fr]{I, 5.3.14}). Comme les
$U_i$ recouvrent $X$ et que $g$ est séparé
(\hyperref[I.5.5.1-fr]{I, 5.5.1, (v)}), $g$ est une immersion de $X$ dans $P$
(\hyperref[I.8.2.8-fr]{I, 8.2.8}). Comme le morphisme de Segre
(\hyperref[II.4.3.3-fr]{4.3.3}) fournit une immersion de $P$ dans un
$\mathbf{P}_k^N$, cela achève de prouver que $X$ est quasi-projectif.

\begin{corollary}[7.4.11]
\label{II.7.4.11-fr}
Une courbe algébrique normale $X$ est isomorphe au schéma induit sur un ouvert
partout dense d'une courbe algébrique normale et complète $\widehat{X}$,
déterminée à un isomorphisme unique près.
\end{corollary}

Si $X_1$, $X_2$ sont deux courbes normales et complètes, il résulte aussitôt
de (\hyperref[II.7.4.9-fr]{7.4.9}) que tout isomorphisme d'un ouvert $U_1$
dense dans $X_1$, sur un ouvert $U_2$ dense dans $X_2$, se prolonge de façon
unique en un isomorphisme de $X_1$ sur $X_2$~; d'où l'assertion d'unicité. Pour
prouver l'existence de $\widehat{X}$, il suffit de remarquer que l'on peut
considérer $X$ comme un sous-schéma d'un fibré projectif $\mathbf{P}_k^n$
(\hyperref[II.7.4.10-fr]{7.4.10}). Soit $\overline{X}$ l'adhérence de $X$ dans
$\mathbf{P}_k^n$ (\hyperref[I.9.5.11-fr]{I, 9.5.11})~; comme $X$ est induit
par $\overline{X}$ sur un ouvert dense dans $\overline{X}$
(\hyperref[I.9.5.10-fr]{I, 9.5.10}), les points génériques $x_i$ des
composantes irréductibles de $X$ sont aussi ceux des composantes irréductibles
de $\overline{X}$, et les $k(x_i)$ sont les mêmes pour ces deux schémas, donc
(\hyperref[II.7.4.1-fr]{7.4.1}) $\overline{X}$ est une courbe algébrique sur
$k$, qui est réduite (\hyperref[I.9.5.9-fr]{I, 9.5.9}) et projective sur $k$
(\hyperref[II.5.5.1-fr]{5.5.1}), donc complète
(\hyperref[II.5.5.3-fr]{5.5.3}). Prenons alors pour $\widehat{X}$ la
\emph{normalisée} de $\overline{X}$, qui est encore complète
(\hyperref[II.7.4.8-fr]{7.4.8})~; en outre, si
$h:\widehat{X}\to\overline{X}$ est le morphisme canonique, la restriction de
$h$ à $h^{-1}(X)$ est un isomorphisme sur $X$ puisque $X$ est normale
(\hyperref[II.6.3.4-fr]{6.3.4}), et comme $h^{-1}(X)$ contient les points
génériques des composantes irréductibles de $\widehat{X}$
(\hyperref[II.6.3.8-fr]{6.3.8}), il est dense dans $\widehat{X}$, ce qui achève
la démonstration.

